\section*{第 \ref{ch:g7:how-animals-breathe} 章——动物如何呼吸}

\begin{solution}{exo:g7:how-animals-breathe:1}
一个大的、薄的、湿润的交换表面，\omterm{prop:g4:journey-of-food:absorb}{血液}（或体液）
在此与外界相接。大：总的穿越面积足够；薄：穿越容易；
湿润：气体溶解后穿过。
\end{solution}

\begin{solution}{exo:g7:how-animals-breathe:2}
肺：\omterm{prop:g3:sorting-living-things:groups}{哺乳动物}、\omterm{prop:g3:sorting-living-things:groups}{鸟类}（\omterm{prop:g4:classifying-animals:classes}{爬行动物}和成年\omterm{prop:g3:sorting-living-things:groups}{两栖动物}也有）。
鳃：\omterm{prop:g3:sorting-living-things:groups}{鱼类}、\omterm{ex:g2:animals-grow-up:frog}{蝌蚪}（河蚌和水生幼虫也有）。
\omterm{def:g7:how-animals-breathe:four}{气管}：蚱蜢、蜜蜂——一般指\omterm{prop:g3:sorting-living-things:groups}{昆虫}。\omterm{def:g7:how-animals-breathe:four}{皮肤
呼吸}：蚯蚓；\omterm{prop:g3:sorting-living-things:groups}{两栖动物}以\omterm{def:g1:the-five-senses:sense}{皮肤}作为辅助。
\end{solution}

\begin{solution}{exo:g7:how-animals-breathe:3}
水从口进入，被推过\omterm{def:g1:animals-around-us:covering}{羽毛}状的鳃，
气体交换即在此处发生——再从鳃盖下方排出：
这是一股单向水流，由口和鳃盖交替
抽动而形成。
\end{solution}

\begin{solution}{exo:g7:how-animals-breathe:4}
通过\omterm{def:g7:how-animals-breathe:four}{气管}：这是一套细小的\omterm{def:g7:how-animals-breathe:four}{气管}网络，
从体表的开口延伸，并\omterm{def:g1:plants-around-us:tree}{分支}到每个器官——空气
被直接送达各处，因此\omterm{prop:g4:journey-of-food:absorb}{血液}需要携带的\omterm{def:g7:what-is-respiration:respiration}{氧气}很少。
\end{solution}

\begin{solution}{exo:g7:how-animals-breathe:5}
离开水后，它\omterm{def:g1:animals-around-us:covering}{羽毛}状的鳃表面塌陷并相互粘连：
大表面几乎缩减为零，
没有表面就无法交换——\omterm{def:g7:what-is-respiration:respiration}{氧气}虽多，
却没有装置来摄取。
\end{solution}

\begin{solution}{exo:g7:how-animals-breathe:6}
\omterm{def:g7:how-animals-breathe:four}{皮肤呼吸}——但湿润的\omterm{def:g6:exploring-our-environment:conditions}{条件}失效：干燥的
\omterm{def:g1:the-five-senses:sense}{皮肤}不能再让气体通过，蚯蚓便在空气中窒息。
大和薄并不能挽救
一个不再湿润的表面。
\end{solution}

\begin{solution}{exo:g7:how-animals-breathe:7}
\omterm{ex:g2:animals-grow-up:frog}{蝌蚪}：用鳃。成体：简单的肺加上湿润的\omterm{def:g1:the-five-senses:sense}{皮肤}——
仅靠\omterm{def:g1:the-five-senses:sense}{皮肤}即可撑过在池底蛰伏的
寒冬。
\end{solution}

\begin{solution}{exo:g7:how-animals-breathe:8}
它的\omterm{def:g7:what-is-respiration:respiration}{呼吸}装置是肺——继承自以\omterm{def:g2:animals-grow-up:born}{乳汁}哺育后代的陆生
祖先——因此无论生活方式多么像鱼，它都必须
浮出水面\omterm{def:g7:what-is-respiration:respiration}{呼吸}空气：若被压入水下便会溺死。外衣
属于水，\omterm{def:g7:what-is-respiration:respiration}{呼吸}装置则属于家族。
\end{solution}

\begin{solution}{exo:g7:how-animals-breathe:9}
携带的空气储备——好比潜水员的气瓶，它藏在翅鞘下方，
其\omterm{def:g7:how-animals-breathe:four}{气管}由此取气。它必须定期返回
水面，以补充气泡中的空气。
\end{solution}

\begin{solution}{exo:g7:how-animals-breathe:10}
(a) 河蚌：水中居民，没有明显的运动，但
鳃位于贝壳内，由泵入的水流冲刷。(b)
蚱蜢：空气中的居民；身体上排列着细小的开口，腹部有节律地伸缩——
以\omterm{def:g7:how-animals-breathe:four}{气管}\omterm{def:g7:what-is-respiration:respiration}{呼吸}。(c) 草蛇：往返于陆地与水域之间的\omterm{prop:g4:classifying-animals:classes}{爬行动物}——
用肺\omterm{def:g7:what-is-respiration:respiration}{呼吸}；游泳时鼻孔朝上，并浮出水面
换气。
\end{solution}

\begin{solution}{exo:g7:how-animals-breathe:11}
\omterm{def:g7:how-animals-breathe:four}{气管}依靠\omterm{def:g1:plants-around-us:tree}{分支}的\omterm{def:g7:how-animals-breathe:four}{气管}输送空气，在毫米尺度上效果极佳，
但在远距离上速度太慢：若身体如鲸鱼一般巨大，
内脏距离任何开口都会太远。这一机制
限定了体型上限，\omterm{prop:g3:sorting-living-things:groups}{昆虫}正生活在这一上限之下。
\end{solution}

\begin{solution}{exo:g7:how-animals-breathe:12}
要求是：身体与外界之间有一个大的、薄的、湿润的接触面，
鳃以\omterm{def:g1:animals-around-us:covering}{羽毛}状的表面满足这一要求，这些表面由水流展开；
肺以数百万个袋状结构满足要求，它们折叠在
湿润的内部空间中；\omterm{def:g7:how-animals-breathe:four}{气管}以最细末端的内表面总和满足要求，
末端保持湿润；\omterm{def:g1:the-five-senses:sense}{皮肤}则以身体整个湿润的外表面满足要求，
相对于较小、纤薄的身体而言，这一面积保持得较大。
四种器官，同一个要求。
\end{solution}

\begin{solution}{pb:g7:how-animals-breathe:1}
\textbf{1.} 鳟鱼苗和刺鱼：用鳃。\omterm{ex:g2:animals-grow-up:frog}{蝌蚪}：
用鳃（长出腿的\omterm{ex:g2:animals-grow-up:frog}{蝌蚪}正处于改造之中）。成年青蛙：用肺
加\omterm{def:g1:the-five-senses:sense}{皮肤}。池塘螺：用肺——往返水面的\omterm{def:g7:what-is-respiration:respiration}{呼吸}者。龙虱：
用\omterm{def:g7:how-animals-breathe:four}{气管}，由其携带的气泡供气。蜉蝣
幼虫：用鳃（\omterm{def:g1:animals-around-us:covering}{羽毛}状的丛状物）。蚯蚓：用\omterm{def:g1:the-five-senses:sense}{皮肤}。

\textbf{2.} 鳃。摆动使水不断更新，让新鲜水
贴着鳃表面流动——在静止的池塘里没有水流，
鳃靠自身摆动制造水流，使含氧的新鲜水
保持在气体交换处。

\textbf{3.} 池塘螺、龙虱、成年青蛙
（以及偶尔到访的草蛇）：都是用肺或气泡\omterm{def:g7:what-is-respiration:respiration}{呼吸}的
\omterm{def:g1:animals-around-us:animal}{动物}——它们从空气中获取\omterm{def:g7:what-is-respiration:respiration}{氧气}，因此
水面就是它们的加油站。

\textbf{4.} \omterm{def:g2:animals-grow-up:metamorphosis}{变态}过程中的\omterm{def:g7:what-is-respiration:respiration}{呼吸}改造：为
无腿\omterm{ex:g2:animals-grow-up:frog}{蝌蚪}服务的鳃正被替换为肺
（以及\omterm{def:g7:what-is-respiration:respiration}{呼吸}用\omterm{def:g1:the-five-senses:sense}{皮肤}），这是即将上岸的幼蛙所需的。

\textbf{5.} 用鳃的\omterm{def:g1:animals-around-us:animal}{动物}最先受到影响：它们的全部
\omterm{def:g7:what-is-respiration:respiration}{氧气}供应来自水中日益减少的溶解氧。往返的
池塘螺、龙虱、青蛙——则取用上方无限的空气，
只是往返得更频繁。

\textbf{6.} 最上层的一毫米与空气接触：\omterm{def:g7:what-is-respiration:respiration}{氧气}
从上方溶入，因此水表薄膜是缺氧池塘中
最不缺氧的水。三刺鱼正
在这最后一处水体前排队。

\textbf{7.} 鳃被堵塞：细泥沉积在
\omterm{def:g1:animals-around-us:covering}{羽毛}状的表面上，阻断气体交换。预计
用鳃\omterm{def:g7:what-is-respiration:respiration}{呼吸}者逃离泥云，而鳃盖泵水
会变得费力，需要痉挛式咳动清除淤泥。

\textbf{8.} \omterm{def:g7:how-animals-breathe:four}{皮肤呼吸}，以及湿润但缺氧的
\omterm{def:g6:exploring-our-environment:components}{环境}：积水的洞穴中\omterm{def:g7:what-is-respiration:respiration}{氧气}太少，
因此蚯蚓浮到表面\omterm{def:g7:what-is-respiration:respiration}{呼吸}，以下方的窒息
换取上方的干燥风险。

\textbf{9.} 按要求填写表格（每种\omterm{def:g1:animals-around-us:animal}{动物}的机制与介质，
“往返者”：池塘螺、龙虱、青蛙）。每列表头
应写明：大而薄且湿润的交换表面，这是
四种机制共同具备的特征。

\textbf{10.} 青蛙：供氧充足，在泥中蛰伏，
以\omterm{def:g7:how-animals-breathe:four}{皮肤呼吸}满足其冬季的少量需求。池螺与龙虱：
处境不佳，体表封闭；它们以休眠
降低需求或依靠所携气泡过冬。\omterm{prop:g3:sorting-living-things:groups}{鱼类}：
只要水中\omterm{def:g7:what-is-respiration:respiration}{氧气}储量尚存即可满足，低温减慢每台机器的运转
（见\cref{ex:g7:what-is-respiration:demand}），鳃
在冰下满足降低后的需求。

\textbf{11.} 蜘蛛的丝钟是固定的储气箱，
而龙虱的气泡是可携带的储气箱：两者都在水下
储存水面空气，两者都在水下
继续服务于一套空气\omterm{def:g7:what-is-respiration:respiration}{呼吸}机制（蜘蛛的书肺与\omterm{def:g7:how-animals-breathe:four}{气管}
装置），却处于错误的介质之中。

\textbf{12.} 池塘在两种介质中提供\omterm{def:g7:what-is-respiration:respiration}{氧气}：
下方稀少且为溶解氧，上方丰富，每种
机器只从一种介质中获取\omterm{def:g7:what-is-respiration:respiration}{氧气}，服务于一类体型的生命：只有
四种机器同时运转，从淤泥到水面薄膜的每一个角落
才都有\omterm{def:g6:exploring-our-environment:components}{生物}栖居。
\end{solution}
